
\section{System Properties and Use Cases}
Multi-agent systems (MAS) increasingly sit in the loop of high-stakes, data-rich applications such as coordination, scheduling, energy, and logistics. They expose multiple attack surfaces at once: valuable \emph{goals} and \emph{private objectives} that an adversary may subvert, abundant \emph{private data} such as user attributes, constraints, and locations, and \emph{communication channels} such as messages, blackboards, and tools. LLM-driven agents amplify both capability and risk: instructions arrive in natural language, tools are invoked via text APIs, and behaviors can be steered or poisoned through subtle prompt-level manipulations. Our aim is to study MAS properties against adversaries in settings that are realistic yet evaluable.


\subsection{Problem Setup}
\label{sec:prob_setup}

We seek MAS problem classes that (i) admit a broad spectrum of attacks (exfiltration, poisoning, spoofing, delay/DoS, collusion), (ii) provide a well-defined, quantitative ground truth for evaluating success and failure, and (iii) are implementable with LLM-based agents.

\medskip
We therefore focus on following system properties:

\paragraph{Agents.} Agents are instantiated by LLMs (optionally tool-augmented), which read natural instructions, exchange messages, and output decisions/actions. This preserves the \emph{agentic interface} used in practice while giving us programmatic control over inputs, channels, and logs.


\pbe{Joint goal.} 
Agents participate in a cooperative task with a well-defined global objective. A shared goal makes adversarial impact \emph{measurable}: subverting coordination (e.g., degraded utility, violated constraints) becomes a scalar signal rather than anecdote. It also enables oracle baselines (optimal or attack-free solutions) for evaluation, i.e., by measuring task utility value gap to oracle, violation counts, regret etc.

\pbe{Agentic data and capabilities.}
Each agent possesses private data, constraints, roles, tools, and actuation. This heterogeneous private state creates a rich privacy and security surface that enables attacks such as exfiltration, impersonation, and capability escalation, while simultaneously exercising the system’s access control and disclosure policies.


\pbe{Inter-agent communication.}
Agents can communicate over addressable channels—pairwise, group, and broadcast—and across modalities such as text, structured JSON, and images. Selective, partially private communication both opens realistic attack vectors, including eavesdropping, spoofing, poisoning, man-in-the-middle, and Sybil coalitions, and enables corresponding defenses, including authentication, redundancy, and audits. Partially private communication can also be needed in practice, e.g., between subsets of agents who belong to the same tenant in enterprise applications. To support realistic experiments and reproducible analysis, the design provides per-message recipient control by the sender, optional encryption and authentication tags, and logged transcripts for ground-truth analysis and post-hoc forensics.
% \eb{Check and update person implementing attacks}


\paragraph{Problem formalization.}
We now introduce a formalization of our cooperative environments, modeled as instruction-augmented DCOPs, see Appendix~\ref{sec:glossary} for full notation glossary. 

\begin{definition}[Instruction-Augmented DCOPs]
\label{def:cpdcop}
An instance is a tuple
\[
\mathcal{P}=\langle A,\ H,\ \eta,\ X,\ D,\ \sigma,\ C,\ \mu,\ o,\ \rho,\ F^\star,\ \Pi\rangle.
\]
\end{definition}

\paragraph{Objective.}
Given a  joint policy $\pi=(\pi_i)_{i\in A}$ and context $c\sim\mu$, for each agent $i$ draw a \emph{joint} assignment over its owned variables
$x_{X_i}\sim \pi_i(\cdot\mid o_i(c))$ independently across $i$
(or set $x_{X_i}=\pi_i(o_i(c))$ if deterministic), and assemble
$x=\big(x_v\big)_{v\in X}$ from the collection $\{x_{X_i}\}_{i\in A}$.
Define the cooperative objective
\[
F^\star(x;c)=\sum_{i\in A}\ \sum_{\beta=1}^{k_i} f_{i,\beta}^\star\big(x_{S_{i,\beta}};c\big),
\]
and the optimization problem
\[
\max_{\ \pi\in\prod_{i\in A}\Pi_i}\;
\mathbb{E}_{c\sim\mu}\ \mathbb{E}_{x\sim \pi(\cdot\mid c)}\big[\,F^\star(x;c)\,\big],
\]
where $x\sim \pi(\cdot\mid c)$ denotes the joint assignment induced by the
agent-wise draws $\{x_{X_i}\}_{i\in A}$. Note that here the $o_i(c)$ is private to agent $i$ and $c$ is not jointly observed by any agents. The pair $(X,F^\star)$ induces a bipartite factor graph $G=(V_\text{var}\cup V_\text{fac},E)$ with $V_\text{var}=X$, $V_\text{fac}=F^\star$, and edge $(x,f)$ iff $x\in \mathrm{scope}(f)$. Algorithms may pass messages on $G$; in practice, we realize it using blackboards.


% \subsection{Extending to more problems.} 
% Our setting supports more complex scenarios that we will keep for future work, however
% we will briefly describe how DCOPs could be extend to other settings:


% \pbe{Competitive and collaborative trade-offs.}

% \pbe{Groups and alliances}

% \pbe{TBD}


\begin{comment}
\section*{Saad's\_Draft: Formal DCOP with Natural Instructions}
\label{sec:context-policy-dcop}


\subsection{Why DCOPs?}
\label{subsec:why-dcop-properties}

\noindent\textbf{Setting.}
We work in the DCOPs of Def.~\ref{def:cpdcop}:
agents $A$ control decision variables $X$ with domains $D$ via $\sigma$;
ground-truth utilities decompose as $F^\star(x;c)=\sum_{i,\beta} f_{i,\beta}^\star(x_{S_{i,\beta}};c)$;
context is drawn from $(C,\mu)$ and observed locally through $o$;
instructions are produced by $\rho(F_i^\star,c)\!\mapsto\! I_i$;
each agent chooses a context-conditioned policy $\pi_i\in\Pi_i$.

\paragraph{Properties of DCOP (with context and instruction).}
\begin{enumerate}
\item \textbf{Cooperative baseline on a fixed tuple.}
Given $\langle A,X,D,\sigma,C,\mu,o,\rho,F^\star,\Pi\rangle$, there is a \emph{single, normative} ground-truth objective
$F^\star(x;c)=\sum_{i,\beta} f_{i,\beta}^\star(x_{S_{i,\beta}};c)$ that all agents jointly optimize.
This makes honest participation well-defined, and any deviation (misreported factor payloads, message/prompt tampering,
topology changes, or policy manipulation) a \emph{measurable anomaly} relative to the same underlying instance.
By contrast, in competitive settings (e.g., general-sum games/POSGs) the “ground truth’’ is ambiguous:
(i) objectives are agent-specific (vector-valued) with no unique welfare aggregator; (ii) multiple equilibria and
equilibrium-selection effects confound attribution; (iii) observed deviations may be rational best responses rather
than attacks; and (iv) private information and incentives entangle signaling with manipulation, making misreporting
indistinguishable from strategic play. The cooperative DCOP baseline avoids these confounds—$\rho(F_i^\star,c)\!\to\!I_i$
separates interface from objective, $F^\star$ fixes the normative target, and attack impact can be isolated and
reproduced on the same tuple.

\item \textbf{Local factorization and causal traceability.}
Small scopes $S_{i,\beta}\!\subseteq\!X$ localize interactions \emph{and} attack surfaces.
By perturbing a factor ($\Delta_{i,\beta}$), an adjacent edge’s timing, or a policy over $X_i$,
we can trace effects through the implicit factor graph at two levels:
(i) \emph{individual}—$\textsf{TargetLift}_i$, per-agent $D_{\mathrm{KL}}$, $\textsf{Leak}(i)$, local latency/messages;
(ii) \emph{population}—$\textsf{CostGap}$, $\textsf{LatencyGain}$, distribution of per-agent KL, spread across cuts/components.
All quantities can be reported per-context $c$ and averaged over $c\!\sim\!\mu$.


\item \textbf{Context-expressive yet analyzable (complexity).}
Context enters explicitly via $(C,\mu)$ and $o$, and we evaluate policies by
$\mathbb{E}_{c\sim\mu}\mathbb{E}_{x\sim\pi(\cdot\mid c)}[F^\star(x;c)]$.
Optimizing a DCOPs is NP-hard in general, but exact methods are exponential
\emph{in the induced width} $w$ of the (implicit) factor graph (e.g., on the
order of $d^{\,w+1}$ for domain size $d$), and anytime/approximate solvers
inherit this structural dependence—yielding analyses parameterized by sparsity
rather than by time horizon. Taking expectations over a finite context set
multiplies running cost by $|C|$ (or by the number of samples), leaving the
worst-case class unchanged. In contrast, sequential competitive models such as
finite-horizon Dec-POMDPs are NEXP-complete, with complexity driven by horizon
length and belief coupling. Our formulation retains expressive context while
avoiding those blow-ups, making attack studies tractable and comparable.


  \item \textbf{Instruction–ground-truth separation for LLM agents.}
  The renderer $\rho(F_i^\star,c)\!\to\!I_i$ exposes only instructions (text/image) and local observations
  to agents while $F^\star$ remains latent. This cleanly models prompt/data poisoning and targeted influence
  as policy-level changes \emph{without} altering the ground-truth environment.



  \item \textbf{Inspectable message passing and instrumentation.}
  Algorithms operate by exchanging structured utilities/bounds/assignments on the implicit graph
  induced by $(X,F^\star)$. This yields precise intervention points (payloads, timing, routing)
  and supports reproducible, stepwise forensics.

  \item \textbf{Graph-theoretic grounding of attacks.}
  Sybil joins, cut-set delays, and multi-spot cost biases correspond to node/edge and cycle operations
  on the induced factor graph, making adversarial budgets and effects analyzable in graph terms.

  \item \textbf{Reproducibility and ablation on the same instance.}
  The tuple $\langle A,X,D,\sigma,C,\mu,o,\rho,F^\star,\Pi\rangle$ fixes all semantics.
  A/B comparisons (\emph{with} vs.\ \emph{without} attack) under identical seeds isolate an attack surface
  (payload vs.\ timing vs.\ instruction) without confounds.



\end{enumerate}

New properties:

private data, private objectives, arbitrary communication patterns.

\begin{enumerate}
    \item roles (e.g., moderator, entity 1, entity 2, entity 3) and capabilities (can send to whom)
    
    \item private communication between subsets of agents
    \begin{itemize}
        \item why do we need it? allows forming groups. Some data is okay (according to CI) to be shared only within a subset of agents (e.g., agents representing users from the same company, e.g., they can decide on priorities based on special release timeline)
        \item what attacks it allows? 
        
        - confidentiality: an outsider can manipulate agents to leak private blackboards it does not have access to. an agent can itself be malicious and leak the blackboard content 
        
        - availability: a malicious agent can start a blackboard with an agent and give it instructions that would delay it.

        - integrity: a malicious agent can give harmful instructions to another agent to override its utility function 

        - influence: a malicious agent can create a colluding ring against another agent 
        
         \end{itemize}
    \item ability to broadcast to all agents 
    \begin{itemize}
        \item why do we need it? some information should be given to everyone. track of global progress. some actions require transparency. enables collaborative problem solving. 

        \item what attacks it allows? a malicious agent can inject the moderator. if there are constraints who can initiate conversations with whom, the global blackboard can override this. 
        
    \end{itemize}
\end{enumerate}

\subsection{Threat Model and Surfaces}
\label{subsec:threat-model-attacks}
\textbf{Adversary goals.} (i) \emph{Confidentiality}: extract private model
parameters, domains, or preferences; (ii) \emph{Availability}: delay/deny
coordination; (iii) \emph{Integrity}: perturb costs/messages/assignments;
(iv) \emph{Influence}: steer the induced policy toward attacker-desired
behaviors.

\noindent\textbf{Capabilities.} (a) Compromised (Byzantine) agent deviating
from protocol; (b) Network adversary that eavesdrops/delays/drops/injects
messages; (c) Colluding set; (d) Sybil identities altering membership/topology;
(e) LLM-surface access to instruction rendering or in-context exemplars.

\noindent\textbf{Surfaces.} (1) Factor payloads and reports
$f_\alpha(\cdot;c)$; (2) Message contents;
(3) Timing/scheduling of messages; (4) Graph topology $G$ (neighbors, degree);
(5) Instructional interfaces (text/image prompts) that condition policies.

\subsection{Global, Context-Aware Evaluation Metrics}
\label{subsec:global-metrics}
We use policy-level, context-aware metrics:
\begin{align}
\textsf{CostGap} &=
\mathbb{E}_{c}\mathbb{E}_{x\sim\pi_{\mathrm{atk}}(\cdot\mid c)}[F(x;c)]
- \mathbb{E}_{c}\mathbb{E}_{x\sim\pi_{0}(\cdot\mid c)}[F(x;c)],
\label{eq:costgap}
\\
\textsf{Stealth} &=
\mathbb{E}_{c}\Big[D_{\mathrm{KL}}\!\big(\pi_{\mathrm{atk}}(\cdot\mid c)\,\Vert\,
\pi_{0}(\cdot\mid c)\big)\Big], \label{eq:stealth} \\
\textsf{TargetLift}(i,a^\star) &=
\mathbb{E}_{c}\Big[\pi_{\mathrm{atk},i}(a^\star\!\mid c) - \pi_{0,i}(a^\star\!\mid c)\Big],
\label{eq:targetlift} \\
\textsf{LatencyGain} &=
\frac{T_{\text{atk}}(\varepsilon)-T_{0}(\varepsilon)}{T_{0}(\varepsilon)}
\quad\text{or}\quad
\frac{K_{\text{atk}}(\varepsilon)-K_{0}(\varepsilon)}{K_{0}(\varepsilon)},
\label{eq:latency} \\
\textsf{Leak}(j) &=
I\!\big(\Theta_j;\,\mathcal{M}\big)
= H(\Theta_j)-H(\Theta_j\mid \mathcal{M}),
\label{eq:leak}
\end{align}
where $T(\varepsilon)$ / $K(\varepsilon)$ are time / rounds to reach an
$\varepsilon$-optimal solution; $\mathcal{M}$ denotes observed messages/timings;
$\Theta_j$ are private parameters of agent $j$ (e.g., coefficients or domains).
\end{comment}



% \subsection{Setup}

% \textbf{Each agent represents a user or some organization with their own private data and objectives.}

% \textbf{Communication is flexible} Anyone can decide to communicate with whoever they like (Sahar: but this actually can be configurable as well?, Mason: Correct, fully configurable) 

% \textbf{Structure is configurable.}


% Communication protocols

% Alignment

% Attacks

% \subsection{What are the properties?}


% Different things that we want to show? 

% soundness, completeness, termination.

% verification of properties?

% Types of games that are implementable? 

% What are properties? 




% \textbf{Travel Planning}

% \textbf{Discussion}

 

